\chapter{Design, Implementation \& Testing}

\section{Overview}

Chapter 3 defined the experimental design and system requirements for the study. This chapter describes how that design was implemented as a modular evaluation pipeline. The pipeline processes original learner sentences, applies prompt variants, stores generated corrections, and evaluates the outputs using grammatical accuracy measures and over-correction risk indicators.

The system developed for this project is not a trainable GEC model. Instead, it is an evaluation pipeline designed to measure how changes in prompt formulation affect correction behaviour. The model parameters remain fixed throughout the experiment, making prompt wording the main variable under investigation.

The implementation was organised into separate stages. Each stage was built as an individual component and tested before being added to the full workflow. This made it easier to identify problems in specific parts of the system, such as dataset formatting, prompt rendering, output storage, or metric calculation.

\textbf{Further implementation details, including hardware, software, library versions, and prompt files, are provided in Appendix~A.}

\section{System Design}

\subsection{Design Rationale}

The system was designed to allow a controlled comparison between prompt conditions. Across all prompt variants, the same learner sentences, reference corrections, model settings, and sentence order were used. This helped ensure that differences in the generated outputs could be linked mainly to prompt design rather than to changes in the data or model configuration.

A within-sentence design was used. In this setup, each learner sentence was processed under every prompt condition. This produced parallel outputs for the same input sentence and allowed direct sentence-level comparison between prompt strategies.

The evaluation followed two separate paths. One path used ERRANT to assess grammatical correction quality. The second path estimated over-correction risk using source-output change and utility-related signals. This separation was needed because a correction may be grammatically acceptable but still change the learner's wording more than necessary.

\subsection{Pipeline Architecture}

Table~\ref{tab:chapter4_pipeline_stages} summarises the workflow implemented in this study.

\begin{table}[htbp]
    \centering
    \caption{Overview of the implemented pipeline stages}
    \label{tab:chapter4_pipeline_stages}
    \begin{tabular}{|p{2cm}|p{4cm}|p{7cm}|}
        \hline
        \textbf{Stage} & \textbf{Component} & \textbf{Purpose} \\
        \hline
        1 & Dataset preparation & Extract the learner sentences and corresponding reference corrections \\
        \hline
        2 & Prompt rendering & Insert each learner sentence into the relevant prompt variant \\
        \hline
        3 & Output generation & Produce a corrected output for each sentence--prompt pair \\
        \hline
        4 & Structured storage & Store each output with its sentence ID, reference correction, and prompt metadata \\
        \hline
        5 & ERRANT evaluation & Compute precision, recall, and $F_{0.5}$ \\
        \hline
        6 & Change-based evaluation & Measure source-output change and utility-related signals \\
        \hline
        7 & OCI computation & Combine source-output change with fluency and meaning-preservation signals \\
        \hline
        8 & Analysis preparation & Generate aggregated files for prompt comparison, sentence-length analysis, robustness checks, trade-off analysis, and qualitative inspection \\
        \hline
    \end{tabular}
\end{table}

\subsection{Data Flow and Storage}

Each output record stored the original learner sentence, the gold reference correction, the generated correction, the sentence identifier, the prompt type, and the prompt version. This format allowed the same outputs to be used for both reference-based evaluation and source-output comparison.

Intermediate outputs were saved in JSONL format. This format was suitable because each sentence--prompt pair could be stored as a separate record. It also made it easier to inspect individual examples, recreate evaluation files, and trace errors back to the original sentence and prompt condition.

The main intermediate files were:

\begin{itemize}
    \item source files containing the original learner sentences;
    \item hypothesis files containing the corrections generated by the model;
    \item reference files containing the gold-standard corrections;
    \item JSONL files storing sentence-level metadata;
    \item CSV files storing metric scores and aggregated results.
\end{itemize}

\section{Implementation Details}

\subsection{Dataset Preparation}

The original learner sentences and gold reference corrections were taken from W\&I + LOCNESS \cite{bryant-etal-2019-bea}. Each sentence was given a stable identifier so that alignment could be maintained across prompt conditions and evaluation files.

Only minimal preprocessing was applied. This was important because the study compares each corrected sentence with the original learner sentence when estimating source-output change and over-correction risk. For this reason, no lemmatisation, paraphrasing, or grammatical normalisation was applied before prompting.

For the sentence-length analysis, token counts were calculated for each source sentence. The sentences were then assigned to short, medium, or long groups using the categories defined in Chapter 3.

\subsection{Prompt Rendering}

The prompt templates were implemented as reusable text structures. Each template included a placeholder for the learner sentence. During execution, this placeholder was replaced with the target sentence to create the full prompt sent to the model.

Four prompt families were included in the experiment:

\begin{itemize}
    \item baseline prompt;
    \item instruction-constrained prompts;
    \item role-based prompts;
    \item few-shot minimal-edit prompts.
\end{itemize}

Several versions were created for the structured prompt families so that prompt refinement could be carried out. Each version was assigned a unique label, such as \texttt{instruction\_v2} or \texttt{role\_v4}. These labels made it possible to group, compare, and analyse outputs by prompt type and version.

Prompt rendering was kept separate from output generation. This allowed the prompt templates to be checked before model execution and ensured that the exact prompt used for any generated output could be reproduced later.

\subsection{Output Generation}

Each learner sentence was processed under every prompt condition. For every generated output, the original sentence, reference correction, prompt type, prompt version, and model correction were stored together.

The same model settings were used for all prompt conditions. The experiments used OpenAI GPT-4o with temperature 0.0, top-p 1.0, a maximum output length of 128 tokens, and seed 1234, as specified in the project configuration file. These settings were kept fixed because changes in decoding parameters could influence correction behaviour and introduce an additional variable. The experiment was intended to compare the effect of prompt design while keeping the sampling configuration constant.

Only light cleaning was applied to the model outputs. This removed formatting artefacts such as unnecessary quotation marks or repeated prompt text. The corrected sentences themselves were not manually edited.

\subsection{ERRANT-Based Evaluation}

Grammatical correction quality was evaluated using ERRANT \cite{bryant-etal-2017-automatic}. For each prompt condition, aligned \texttt{.src}, \texttt{.hyp}, and \texttt{.ref} files were created. These represented the original learner sentence, the model output, and the gold reference correction.

ERRANT was used to calculate precision, recall, and $F_{0.5}$. The $F_{0.5}$ score was treated as the main grammatical accuracy metric in this study because it gives more weight to precision. This is relevant for GEC, where unnecessary edits can reduce the quality of a correction even if the output is fluent.

\subsection{Change-Based Evaluation}

A separate evaluation stage calculated source-output change and utility-related signals for each model output. This stage compared the source sentence directly with the model output, rather than comparing the output with the gold reference correction.

The following metrics were computed:

\begin{table}[htbp]
    \centering
    \caption{Change-based and utility-related metrics used for over-correction risk analysis}
    \label{tab:style_metrics}
    \begin{tabular}{|p{5cm}|p{8cm}|}
        \hline
        \textbf{Metric} & \textbf{Purpose} \\
        \hline
        Word-level Levenshtein distance & Measures the amount of surface rewriting \\
        \hline
        Edit density & Normalises edit distance by sentence length \\
        \hline
        $\Delta TTR$ & Measures change in lexical diversity \\
        \hline
        $\Delta R$ & Measures change in readability \\
        \hline
        Cosine similarity & Estimates source-output similarity and is used as a meaning-preservation signal \\
        \hline
        $\Delta$ fluency & Measures whether the corrected output improves surface fluency relative to the source \\
        \hline
    \end{tabular}
\end{table}

Sentence-level fluency change was calculated as the difference between the fluency score assigned to the corrected output and the fluency score assigned to the original source sentence.

These metrics do not prove on their own that an edit is unnecessary. Instead, they provide measurable signals about how much the text has changed, whether meaning appears to be preserved, and whether fluency has improved. When interpreted together with ERRANT results and qualitative examples, they help indicate whether a prompt tends to preserve the learner's wording or rewrite it more heavily.

\subsection{OCI Calculation}

In this study, the Composite Over-Correction Index (OCI) is used as an indicator of over-correction risk. It combines source-output change with a utility adjustment. The purpose of the utility adjustment is to avoid treating every change as equally problematic: a correction that changes the source sentence is less likely to be considered risky if it improves fluency and preserves meaning.

A source-output change score was first calculated using several change-based signals. Each component was normalised with min--max scaling:

\begin{equation}
\text{norm}(x) = \frac{x - \min(x)}{\max(x) - \min(x)}
\end{equation}

\noindent The source-output change component was calculated as:

\begin{equation}
OCI_{\text{change}} = w_1 \cdot \text{norm}(D_{\text{edit}}) 
+ w_2 \cdot \text{norm}(E_{\text{density}}) 
+ w_3 \cdot \text{norm}(\Delta TTR) 
+ w_4 \cdot \text{norm}(\Delta R) 
+ w_5 \cdot (1 - \text{CosSim})
\end{equation}

The weights were set before analysis:

\[
w_1 = 0.35,\quad
w_2 = 0.15,\quad
w_3 = 0.20,\quad
w_4 = 0.15,\quad
w_5 = 0.15
\]

Edit distance was given the largest weight because direct textual change is the clearest signal of rewriting. Edit density was included to account for sentence length, while lexical diversity, readability, and cosine similarity captured broader changes in vocabulary, structure, and overall similarity.

The utility component was then calculated as:

\begin{equation}
Utility = 0.5 \cdot \text{norm}(\Delta_{\text{fluency}}) + 0.5 \cdot \text{CosSim}
\end{equation}

Here, $\Delta_{\text{fluency}}$ represents the change in fluency between the source sentence and the corrected output, while $\text{CosSim}$ represents source-output similarity and is used as a meaning-preservation signal. Cosine similarity is used in two directions: $(1-\text{CosSim})$ contributes to the source-output change score, while $\text{CosSim}$ contributes to the utility adjustment. This means that low similarity increases the estimated change, whereas high similarity reduces the final risk score when meaning appears to be preserved. The two utility terms were weighted equally so that fluency improvement and meaning preservation contributed equally to the utility adjustment.

The final OCI score was calculated as:

\begin{equation}
OCI = OCI_{\text{change}} \cdot (1 - Utility)
\end{equation}

Both the utility term and the final OCI value were clipped to the range $[0,1]$. With this scoring approach, a correction receives a higher over-correction risk score when it creates greater distance from the source sentence without improving fluency or preserving meaning. In contrast, a correction can involve visible surface changes but still receive a lower OCI value if those changes improve fluency and keep the original meaning intact.

The role of the utility adjustment is therefore to reduce the effect of source-output change when that change appears useful. It is not intended to decide by itself whether an edit is necessary. OCI was calculated first at sentence level. The sentence-level scores were then grouped and compared by prompt variant, prompt strategy, and sentence-length category. Because OCI combines several indirect signals, it is used for comparing prompt conditions rather than for making absolute judgements about individual corrections.

\section{Analysis Outputs}

The pipeline produced the metric files and aggregated results required for the analyses reported in Chapter 5. These outputs supported prompt refinement, final prompt strategy comparison, sentence-length analysis, OCI robustness testing, trade-off analysis, and qualitative inspection.

For each analysis, sentence-level records were retained so that aggregate results could be traced back to individual examples. This was important because the project evaluates both correction accuracy and over-correction risk. Aggregate scores such as $F_{0.5}$ and OCI were therefore interpreted alongside sentence-level outputs during the results and discussion stage.

\section{Testing and Validation}

\subsection{Data Validation}

Dataset alignment was checked before evaluation. The number of lines in the source, hypothesis, and reference files was verified so that each model output matched the correct source and reference sentence.

Sentence identifiers were also checked for uniqueness and consistency across prompt outputs. This was necessary because incorrect alignment would make both ERRANT scoring and sentence-level OCI comparison unreliable.

\subsection{Pipeline Testing}

Each stage of the pipeline was tested on a small subset of sentences before the full dataset was processed. This included prompt rendering, output storage, ERRANT file generation, change-based metric calculation, and OCI aggregation.

Running the pipeline on a smaller subset first made it easier to identify formatting errors, path issues, and unexpected model outputs before carrying out the full experiment.

\subsection{Metric Verification}

Metric outputs were manually inspected using selected examples. Minimal corrections were expected to produce low edit distance and low OCI values. Heavily rewritten outputs were expected to produce higher OCI values, unless the source-output change was offset by fluency gain and meaning preservation.

Cosine similarity was also checked to make sure that its direction matched the two parts of the OCI formula. It was inverted as $(1 - \text{CosSim})$ in the source-output change component, where higher values indicate greater difference. It was used directly in the utility component, where higher values indicate stronger meaning preservation.

\subsection{Prompt Ranking Validation}

Prompt ranking was checked to ensure that prompt selection was not based only on grammatical accuracy. Both $F_{0.5}$ and OCI were considered when identifying the strongest prompt variants.

This validation step was needed because the project examines over-correction risk as well as correction accuracy. A prompt with the highest $F_{0.5}$ score was therefore not automatically treated as the best overall prompt if it also produced substantially higher over-correction risk, as indicated by OCI.

\subsection{Reproducibility and Error Handling}

Reproducibility was supported through fixed prompts, fixed model settings, saved intermediate files, and script-based evaluation. File naming conventions were kept consistent across experiments so that each output could be traced back to its prompt condition.

Basic safeguards were included to reduce the risk of silent failures. These included checking file paths, enforcing text encoding, validating output counts, and inspecting empty or malformed model outputs before evaluation.